% open-logic.tex
% compiles to produce complete text using open-logic-dev style

\documentclass[../include/open-logic-part]{subfiles}

\begin{document}

\clearpage

\begin{editorial}
Dit bestand laadt alle inhoud die in het Open Logic Project is opgenomen.
Wanneer redactionele opmerkingen zoals deze zichtbaar zijn, betekent dit dat
het bestand is gecompileerd zonder rekening te houden met de manier waarop het
materiaal wordt gepresenteerd. Als u dit kunt lezen, is het waarschijnlijk
\emph{niet raadzaam} om met deze pdf les te geven of eruit te studeren.

Het Open Logic Project biedt veel mogelijkheden om een tekst samen te
stellen die beter geschikt is voor onderwijs of zelfstudie. Standaard worden
bijvoorbeeld alle logische operatoren als primitief behandeld en worden in
bewijzen alle gevallen voor alle operatoren uitgewerkt. Het is echter veel beter
om enkele van die gevallen als opgaven over te laten. Ook is het Open Logic
Project nog in ontwikkeling. Om samenwerking en verbetering te bevorderen,
wordt materiaal ook opgenomen als het alleen nog in conceptvorm is of als er nog
opgaven in ontbreken, enzovoort. In een pdf die voor een cursus wordt
samengesteld, worden zulke paragrafen weggelaten.

Bekijk voor pdf's die beter geschikt zijn voor onderwijs en studie de
\href{http://builds.openlogicproject.org/}{voorbeeldcursussen op de
  OLP-website}. Om zelf een pdf samen te stellen, kunt u beginnen met het
\href{https://github.com/OpenLogicProject/OpenLogic/tree/master/courses/sample}{voorbeelddriverbestand}
of de bronbestanden van de afgeleide leerboeken raadplegen
voor uitgebreidere en geavanceerdere voorbeelden.
\end{editorial}

\olimport[sets-functions-relations]{sets-functions-relations-complete}

\olimport[propositional-logic]{propositional-logic}

\olimport[first-order-logic]{first-order-logic}

\olimport[model-theory]{model-theory}

\olimport[computability]{computability}

\olimport[turing-machines]{turing-machines}

\olimport[incompleteness]{incompleteness}

\olimport[second-order-logic]{second-order-logic}

\olimport[lambda-calculus]{lambda-calculus}

\olimport[many-valued-logic]{many-valued-logic}

\olimport[normal-modal-logic]{normal-modal-logic}

\olimport[applied-modal-logic]{applied-modal-logic}

\olimport[intuitionistic-logic]{intuitionistic-logic}

\olimport[counterfactuals]{counterfactuals}

\olimport[set-theory]{set-theory}

\olimport[methods]{methods}

\olimport[history]{history}

\olimport[reference]{reference}
\end{document}

